\chapter{Concorrência Revisitada}
\label{CH:LOCK2}

Obter simultaneamente um bom desempenho paralelo, correção apesar da concorrência e código compreensível é um grande desafio no design de núcleos. O uso direto de bloqueios é o melhor caminho para 
a correção, mas nem sempre é possível. Este capítulo destaca exemplos em que o xv6 é forçado a usar bloqueios de forma complexa, e exemplos onde o xv6 utiliza técnicas semelhantes a bloqueios, mas 
não bloqueios propriamente ditos.

\section{Padrões de Bloqueio}

Itens em cache são frequentemente um desafio para bloquear. Por exemplo, o cache de blocos do sistema de arquivos \lineref{kernel/bio.c:/^struct/} armazena cópias de até {\tt NBUF} blocos de disco. 
É vital que um dado bloco de disco tenha no máximo uma cópia no cache; caso contrário, diferentes processos podem fazer alterações conflitantes em diferentes cópias do que deveria ser o mesmo 
bloco. Cada bloco em cache é armazenado em uma {\tt struct buf} \lineref{kernel/buf.h:1}. Uma {\tt struct buf} possui um campo de bloqueio que ajuda a garantir que apenas um processo use um dado 
bloco de disco por vez. No entanto, esse bloqueio não é suficiente: e se um bloco não estiver presente no cache, e dois processos quiserem usá-lo ao mesmo tempo? Não há {\tt struct buf} (uma vez 
que o bloco ainda não está em cache), e, portanto, não há nada para bloquear. O xv6 lida com essa situação associando um bloqueio adicional ({\tt bcache.lock}) ao conjunto de identidades dos blocos 
em cache. O código que precisa verificar se um bloco está em cache (por exemplo, {\tt bget} \lineref{kernel/bio.c:/^bget/}), ou alterar o conjunto de blocos em cache, deve manter {\tt bcache.lock}; 
após esse código ter encontrado o bloco e a {\tt struct buf} que precisa, ele pode liberar {\tt bcache.lock} e bloquear apenas o bloco específico. Este é um padrão comum: um bloqueio para o 
conjunto de itens, além de um bloqueio por item.

Normalmente, a mesma função que adquire um bloqueio o liberará. Mas uma forma mais precisa de ver as coisas é que um bloqueio é adquirido no início de uma sequência que deve parecer atômica, e 
liberado quando essa sequência termina. Se a sequência começa e termina em funções diferentes, ou em threads diferentes, ou em CPUs diferentes, então a aquisição e liberação do bloqueio devem fazer 
o mesmo. A função do bloqueio é forçar outros usos a esperar, não fixar um pedaço de dado a um agente particular. Um exemplo é a {\tt acquire} em {\tt yield} \lineref{kernel/proc.c:/^yield/}, que é 
liberada na thread do escalonador, em vez de no processo adquirente. Outro exemplo é a {\tt acquiresleep} em {\tt ilock} \lineref{kernel/fs.c:/^ilock/}; esse código frequentemente dorme enquanto lê 
do disco; pode acordar em uma CPU diferente, o que significa que o bloqueio pode ser adquirido e liberado em CPUs diferentes.

Liberar um objeto que é protegido por um bloqueio embutido no objeto é um negócio delicado, uma vez que possuir o bloqueio não é suficiente para garantir que a liberação seria correta. O caso 
problemático surge quando algum outro thread está esperando em {\tt acquire} para usar o objeto; liberar o objeto implicitamente libera o bloqueio embutido, o que fará com que o thread em espera 
funcione de maneira inadequada. Uma solução é rastrear quantas referências ao objeto existem, para que ele só seja liberado quando a última referência desaparecer. Veja {\tt pipeclose} 
\lineref{kernel/pipe.c:/^pipeclose/} para um exemplo; {\tt pi->readopen} e {\tt pi->writeopen} rastreiam se o pipe tem descritores de arquivo que se referem a ele.

Geralmente, vê-se bloqueios em torno de sequências de leituras e gravações em conjuntos de itens relacionados; os bloqueios garantem que outros threads vejam apenas sequências concluídas de 
atualizações (desde que eles também bloqueiem). E quanto a situações em que a atualização é uma simples gravação em uma única variável compartilhada? Por exemplo, 
\texttt{setkilled} e \texttt{killed} \lineref{kernel/proc.c:/^setkilled/} bloqueiam em torno de seus usos simples de
\lstinline{p->killed}. Se não houvesse bloqueio, um thread poderia gravar 
\lstinline{p->killed} ao mesmo tempo que outro thread o lê. Isso é uma \indextextx{competição}, e a especificação da linguagem C diz que uma competição resulta em \indextext{comportamento 
indefinido}, o que significa que o programa pode falhar ou produzir resultados incorretos\footnote{``Threads and data races'' em \url{
https://en.cppreference.com/w/c/language/memory_model}}. Os bloqueios previnem a competição e evitam o comportamento indefinido.

Uma razão pela qual as competições podem quebrar programas é que, se não houver bloqueios ou construções equivalentes, o compilador pode gerar código de máquina que lê e grava memória de maneiras 
bastante diferentes do código C original. Por exemplo, o código de máquina de um thread chamando
\texttt{killed} poderia copiar \lstinline{p->killed} para um registrador e
ler apenas esse valor em cache; isso significaria que o thread
poderia nunca ver quaisquer gravações em
\lstinline{p->killed}. Os bloqueios previnem tal cache.

% sleep locks.
% bloqueios de hand-over-hand no namei.
% uso de refcount pelo namei para evitar alterações indesejadas,
% e bloqueio ao usá-lo.
% exemplo de onde o deadlock é evitado? namei?
% gerar uma thread para evitar um problema de ordem de bloqueio?

\section{Padrões Semelhantes a Bloqueios}

Em muitos lugares, o xv6 usa uma contagem de referências ou um flag de maneira semelhante a um bloqueio para indicar que um objeto está alocado e não deve ser liberado ou reutilizado. O {\tt p-
>state} de um processo atua dessa forma, assim como as contagens de referências nas estruturas {\tt file}, {\tt inode} e {\tt buf}. Embora em cada caso um bloqueio proteja o flag ou a contagem de 
referências, é esta última que impede que o objeto seja liberado prematuramente.

O sistema de arquivos usa contagens de referência de {\tt struct inode} como uma espécie de bloqueio compartilhado que pode ser mantido por múltiplos processos, a fim de evitar deadlocks que 
ocorreriam se o código usasse bloqueios comuns. Por exemplo, o loop em {\tt namex} \lineref{kernel/fs.c:/^namex/} bloqueia o diretório nomeado por cada componente do caminho em sequência. No 
entanto, {\tt namex} deve liberar cada bloqueio no final do loop, pois se mantiver múltiplos bloqueios, pode ocorrer um deadlock consigo mesmo se o caminho incluir um ponto (por exemplo, {\tt 
a/./b}). Também pode haver deadlock com uma busca concorrente envolvendo o diretório e {\tt ..}. Como explicado no Capítulo~\ref{CH:FS}, a solução é que o loop leve o inode do diretório para a 
próxima iteração com sua contagem de referência incrementada, mas não bloqueada.

Alguns itens de dados são protegidos por diferentes mecanismos em momentos diferentes e podem, às vezes, ser protegidos de acesso concorrente implicitamente pela estrutura do código do xv6, em vez 
de por bloqueios explícitos. Por exemplo, quando uma página física está livre, ela é protegida por \texttt{kmem.lock} \lineref{kernel/kalloc.c:/^. kmem;/}. Se a página for então alocada como um 
pipe \lineref{kernel/pipe.c:/^pipealloc/}, ela é protegida por um bloqueio diferente (o embutido \lstinline{pi->lock}). Se a página for realocada para a memória do usuário de um novo processo, ela 
não é protegida por um bloqueio. Em vez disso, o fato de que o alocador não dará essa página a nenhum outro processo (até que seja liberada) a protege de acesso concorrente. A propriedade da 
memória de um novo processo é complexa: primeiro, o pai aloca e manipula-a em {\tt fork}, depois o filho a utiliza, e (após a saída do filho) o pai novamente possui a memória e a passa para {\tt 
kfree}. Aqui estão duas lições: um objeto de dados pode ser protegido de concorrência de maneiras diferentes em diferentes pontos de seu ciclo de vida, e a proteção pode assumir a forma de 
estrutura implícita em vez de bloqueios explícitos.

Um exemplo final semelhante a bloqueios é a necessidade de desabilitar interrupções em torno de chamadas para {\tt mycpu()} \lineref{kernel/proc.c:/^myproc/}. Desabilitar interrupções faz com que o 
código chamador seja atômico em relação a interrupções de timer que poderiam forçar uma troca de contexto, e assim mover o processo para uma CPU diferente.

\section{Sem bloqueios de forma alguma}

Existem alguns lugares onde o xv6 compartilha dados mutáveis sem bloqueios. Um deles está na implementação de spinlocks, embora possa-se ver as instruções atômicas do RISC-V como dependendo de 
bloqueios implementados em hardware. Outro é a variável {\tt started} em {\tt main.c} \lineref{kernel/main.c:/^volatile/}, usada para evitar que outras CPUs sejam executadas até que a CPU zero 
tenha terminado de inicializar o xv6; o {\tt volatile} garante que o compilador realmente gere instruções de carregamento e armazenamento.

O xv6 contém casos em que uma CPU ou thread escreve alguns dados, e outra CPU ou thread lê os dados, mas não há um bloqueio específico dedicado a proteger esses dados. Por exemplo, em {\tt fork}, o 
pai escreve as páginas de memória do usuário do filho, e o filho (um thread diferente, talvez em uma CPU diferente) lê essas páginas; nenhum bloqueio protege explicitamente essas páginas. Isso não 
é estritamente um problema de bloqueio, uma vez que o filho não começa a executar até que o pai tenha terminado de escrever. É um potencial problema de ordenação de memória (veja o 
Capítulo~\ref{CH:LOCK}), pois, sem uma barreira de memória, não há razão para esperar que uma CPU veja as gravações de outra CPU. No entanto, como o pai libera bloqueios e o filho adquire bloqueios 
ao iniciar, as barreiras de memória em {\tt acquire} e {\tt release} garantem que a CPU do filho veja as gravações do pai.

\section{Padrões Semelhantes a Bloqueios}

Em muitos lugares, o xv6 usa uma contagem de referências ou um flag de maneira semelhante a um bloqueio para indicar que um objeto está alocado e não deve ser liberado ou reutilizado. O {\tt p-
>state} de um processo atua dessa forma, assim como as contagens de referências nas estruturas {\tt file}, {\tt inode} e {\tt buf}. Embora em cada caso um bloqueio proteja o flag ou a contagem de 
referências, é esta última que impede que o objeto seja liberado prematuramente.

O sistema de arquivos usa contagens de referência de {\tt struct inode} como uma espécie de bloqueio compartilhado que pode ser mantido por múltiplos processos, a fim de evitar deadlocks que 
ocorreriam se o código usasse bloqueios comuns. Por exemplo, o loop em {\tt namex} \lineref{kernel/fs.c:/^namex/} bloqueia o diretório nomeado por cada componente do caminho em sequência. No 
entanto, {\tt namex} deve liberar cada bloqueio no final do loop, pois se mantiver múltiplos bloqueios, pode ocorrer um deadlock consigo mesmo se o caminho incluir um ponto (por exemplo, {\tt 
a/./b}). Também pode haver deadlock com uma busca concorrente envolvendo o diretório e {\tt ..}. Como explicado no Capítulo~\ref{CH:FS}, a solução é que o loop leve o inode do diretório para a 
próxima iteração com sua contagem de referência incrementada, mas não bloqueada.

Alguns itens de dados são protegidos por diferentes mecanismos em momentos diferentes e podem, às vezes, ser protegidos de acesso concorrente implicitamente pela estrutura do código do xv6, em vez 
de por bloqueios explícitos. Por exemplo, quando uma página física está livre, ela é protegida por \texttt{kmem.lock} \lineref{kernel/kalloc.c:/^. kmem;/}. Se a página for então alocada como um 
pipe \lineref{kernel/pipe.c:/^pipealloc/}, ela é protegida por um bloqueio diferente (o embutido \lstinline{pi->lock}). Se a página for realocada para a memória do usuário de um novo processo, ela 
não é protegida por um bloqueio. Em vez disso, o fato de que o alocador não dará essa página a nenhum outro processo (até que seja liberada) a protege de acesso concorrente. A propriedade da 
memória de um novo processo é complexa: primeiro, o pai aloca e manipula-a em {\tt fork}, depois o filho a utiliza, e (após a saída do filho) o pai novamente possui a memória e a passa para {\tt 
kfree}. Aqui estão duas lições: um objeto de dados pode ser protegido de concorrência de maneiras diferentes em diferentes pontos de seu ciclo de vida, e a proteção pode assumir a forma de 
estrutura implícita em vez de bloqueios explícitos.

Um exemplo final semelhante a bloqueios é a necessidade de desabilitar interrupções em torno de chamadas para {\tt mycpu()} \lineref{kernel/proc.c:/^myproc/}. Desabilitar interrupções faz com que o 
código chamador seja atômico em relação a interrupções de timer que poderiam forçar uma troca de contexto, e assim mover o processo para uma CPU diferente.

\section{Sem bloqueios de forma alguma}

Existem alguns lugares onde o xv6 compartilha dados mutáveis sem bloqueios. Um deles está na implementação de spinlocks, embora possa-se ver as instruções atômicas do RISC-V como dependendo de 
bloqueios implementados em hardware. Outro é a variável {\tt started} em {\tt main.c} \lineref{kernel/main.c:/^volatile/}, usada para evitar que outras CPUs sejam executadas até que a CPU zero 
tenha terminado de inicializar o xv6; o {\tt volatile} garante que o compilador realmente gere instruções de carregamento e armazenamento.

O xv6 contém casos em que uma CPU ou thread escreve alguns dados, e outra CPU ou thread lê os dados, mas não há um bloqueio específico dedicado a proteger esses dados. Por exemplo, em {\tt fork}, o 
pai escreve as páginas de memória do usuário do filho, e o filho (um thread diferente, talvez em uma CPU diferente) lê essas páginas; nenhum bloqueio protege explicitamente essas páginas. Isso não 
é estritamente um problema de bloqueio, uma vez que o filho não começa a executar até que o pai tenha terminado de escrever. É um potencial problema de ordenação de memória (veja o 
Capítulo~\ref{CH:LOCK}), pois, sem uma barreira de memória, não há razão para esperar que uma CPU veja as gravações de outra CPU. No entanto, como o pai libera bloqueios e o filho adquire bloqueios 
ao iniciar, as barreiras de memória em {\tt acquire} e {\tt release} garantem que a CPU do filho veja as gravações do pai.

\section{Paralelismo}

O bloqueio é principalmente sobre suprimir o paralelismo em prol da
correção. Como o desempenho também é importante, os designers de núcleos
frequentemente precisam pensar em como usar bloqueios de uma maneira que
alcance tanto a correção quanto permita o paralelismo. Embora o xv6
não seja sistematicamente projetado para alto desempenho, ainda vale a pena
considerar quais operações do xv6 podem ser executadas em paralelo, e quais
podem entrar em conflito nos bloqueios.

Os pipes no xv6 são um exemplo de paralelismo relativamente bom. Cada pipe tem
seu próprio bloqueio, de modo que diferentes processos podem ler e escrever
em diferentes pipes em paralelo em diferentes CPUs. Para um dado pipe,
no entanto, o escritor e o leitor devem esperar um pelo outro para liberar o
bloqueio; eles não podem ler/gravar o mesmo pipe ao mesmo tempo. Também é
verdade que uma leitura de um pipe vazio (ou uma gravação em um pipe cheio)
deve bloquear, mas isso não se deve ao esquema de bloqueio.

A troca de contexto é um exemplo mais complexo. Dois threads do núcleo, cada um
executando em sua própria CPU, podem chamar {\tt yield}, {\tt sched} e {\tt
  swtch} ao mesmo tempo, e as chamadas serão executadas em paralelo. Os
threads mantêm cada um um bloqueio, mas são bloqueios diferentes, então não
precisam esperar um pelo outro. Uma vez em {\tt scheduler}, no entanto, as duas
CPUs podem entrar em conflito nos bloqueios enquanto buscam na tabela de processos
por um que esteja {\tt RUNNABLE}. Ou seja, o xv6 provavelmente obterá um
benefício de desempenho de múltiplas CPUs durante a troca de contexto, mas
talvez não tanto quanto poderia.

Outro exemplo é as chamadas concorrentes para {\tt fork} de diferentes
processos em diferentes CPUs. As chamadas podem ter que esperar uma pela outra
pelo {\tt pid\_lock} e {\tt kmem.lock}, e pelos bloqueios por processo
necessários para buscar na tabela de processos por um processo {\tt UNUSED}. Por
outro lado, os dois processos que estão fazendo fork podem copiar as páginas de
memória do usuário e formatar as páginas da tabela de páginas completamente em
paralelo.

O esquema de bloqueio em cada um dos exemplos acima sacrifica o desempenho
paralelo em certos casos. Em cada caso, é possível obter mais paralelismo
usando um design mais elaborado. Se isso vale a pena depende de detalhes:
com que frequência as operações relevantes são invocadas, quanto tempo o código
passa com um bloqueio disputado mantido, quantas CPUs podem estar executando
operações conflitantes ao mesmo tempo, se outras partes do código são gargalos
mais restritivos. Pode ser difícil adivinhar se um determinado esquema de
bloqueio pode causar problemas de desempenho, ou se um novo design é
significativamente melhor, então a medição em cargas de trabalho realistas é
frequentemente necessária.

\section{Exercícios}

\begin{enumerate}

\item Modifique a implementação do pipe do xv6 para permitir que uma leitura e
  uma gravação no mesmo pipe prossigam em paralelo em diferentes CPUs.

\item Modifique o \texttt{scheduler()} do xv6 para reduzir a contenção de
  bloqueios quando diferentes CPUs estão procurando por processos executáveis ao
  mesmo tempo.

\item Elimine algumas das serializações no \texttt{fork()} do xv6.

\end{enumerate}
